\documentclass[11pt]{article}
\usepackage[a4paper,margin=20mm]{geometry}
\usepackage{amsmath,amssymb,amsthm,microtype}
\usepackage[colorlinks=true,linkcolor=blue,urlcolor=blue]{hyperref}
\newtheorem{theorem}{Theorem}
\newtheorem{corollary}[theorem]{Corollary}
\title{Joint polynomial circuits for permutation and monomial unitaries}
\author{TNLean contributors}
\date{3 October 2026}
\begin{document}
\maketitle

Consecutive operator Schmidt rank bounds admit an exact circuit with
polynomial gate count in the chain length when the maximal rank is fixed.
The general construction has a rank-dependent exponent. For permutation
unitaries, the following argument instead gives polynomial dependence
jointly on the chain length and the maximal rank. The probability gap of
a permutation permits an exact reversible computation from finite-precision
canonical tensors. Combining this computation with diagonal phase evaluation
also gives joint polynomial operator-norm approximation of monomial unitaries.
The mathematical results below are not asserted to be formalized in Lean.

\section{Joint polynomial circuits for permutation and monomial MPUs}

The following argument concerns nonuniform circuit existence.
The local dimension $d\ge2$ is fixed. Gates are arbitrary one-
and two-qudit unitaries, or only adjacent two-qudit unitaries after
the routing described below. Every work register is initialized
and returned to the computational basis vector $|0\rangle$.
The argument does not assert efficient conversion of a supplied
tensor representation into canonical tensors or into the resulting
gate parameters. It is a mathematical proof, not a Lean
formalization.

\begin{theorem}\label{thm:joint-polynomial-permutation}
Let $N,D\ge1$, and let $\pi$ be a permutation of
$(\mathbb Z/d\mathbb Z)^N$. Suppose that its permutation unitary
$P_\pi|y\rangle=|\pi(y)\rangle$ has operator Schmidt rank at most
$D$ across every consecutive cut.
There is an exact circuit $C$ for $P_\pi$ with gate count and
auxiliary space polynomial jointly in $N,D$, with exponents
depending only on $d$. In particular, for a constant $c_d>0$, put
\[
 B=c_d\bigl(N+\log_2(N+1)+\log_2(D+1)+1\bigr),\qquad
 S=c_dN^2(D+1)^4B^2.
\]
The constants may be chosen so that $O_d(S)$ arbitrary one- and
two-qudit gates and $O_d(S)$ auxiliary qudits suffice.
On a line, $O_d(S^2)$ adjacent two-qudit gates suffice with the same
auxiliary bound. The exact implementation identity is
$CJ=JP_\pi$, including the scalar phase, where
$J\psi=\psi\otimes|0\rangle$ initializes all auxiliaries.
\end{theorem}

\begin{proof}
Normalize the coefficient vector of the permutation unitary:
\[
 \psi(a,y)=d^{-N/2}(P_\pi)_{a,y}.
\]
Its norm is one. Pairing $a_j,y_j$ site by site identifies it with
a vector on $N$ sites of local dimension $d^2$. Its Schmidt ranks
are the operator Schmidt ranks of $P_\pi$, hence are at most $D$.
Successive Schmidt decompositions give a right-canonical
open-boundary MPS
\[
 \psi(a,y)=A_1(a_1,y_1)\cdots A_N(a_N,y_N),\qquad
 A_j(a,b):E_j\longrightarrow E_{j-1},
\]
where $\dim E_j\le D$, $E_0=E_N=\mathbb C$, and
\[
 \sum_{a,b=0}^{d-1} A_j(a,b)A_j(a,b)^\dagger=I_{E_{j-1}}.
\]
Every individual matrix has operator norm at most one.
This canonical form is the same bounded-contraction ingredient
used in the diagonal approximation argument of
\nolinkurl{docs/audits/2026-10-02_mpu_diagonal_circuits.tex},
Section~3, and in the minimal-representation
step of Section~5 of the exact MPU circuit note.

\subsection*{Output-prefix probabilities}

For a fixed input word $y$ and output prefix
$z=(z_1,\ldots,z_k)$, define
\[
 p(z\mid y)=\sum_{a_{k+1},\ldots,a_N}
       |\psi(z_1,\ldots,z_k,a_{k+1},\ldots,a_N;y)|^2.
\]
Because $\pi$ has exactly one output on each input,
\[
 p(z\mid y)=
 \begin{cases}
  d^{-N},& z\text{ is the prefix of }\pi(y),\\
  0,&\text{otherwise}.
 \end{cases}
\]
The probability is an MPS contraction. Start with the scalar
$R_N=1$ and, in order $j=N,\ldots,1$, apply
\[
 R_{j-1}=\mathcal T_j(R_j),\qquad
 \mathcal T_j(X)=
 \begin{cases}
  A_j(z_j,y_j)X A_j(z_j,y_j)^\dagger,&j\le k,\\[2pt]
  \displaystyle\sum_{a=0}^{d-1}
       A_j(a,y_j)X A_j(a,y_j)^\dagger,&j>k.
 \end{cases}
\]
The final scalar $R_0$ is $p(z\mid y)$.
Every $\mathcal T_j$ is positive and subunital:
$0\le\mathcal T_j(I)\le I$. Therefore $0\le R_j\le I$.
It is also a contraction of the operator norm on Hermitian
matrices. Indeed, $-\|X\|I\le X\le\|X\|I$ implies
$-\|X\|I\le\mathcal T_j(X)\le\|X\|I$.
This is the reason the probability computation does not require
a lower bound on any Schmidt value.

\subsection*{A uniform finite-precision estimate}

Let $u=2^{-b}$, and round every real and imaginary tensor entry
to an integer multiple of $u$, with absolute error at most $u$.
Write the rounded matrices as $\widehat A_j$. Then
\[
 \|\widehat A_j-A_j\|\le2Du.
\]
For Hermitian $X$, define $\widehat{\mathcal T}_j$ using the
rounded matrices but exact arithmetic. If $Du\le1$, the difference
of the two transfer maps satisfies
\[
 \|(\widehat{\mathcal T}_j-\mathcal T_j)(X)\|
 \le c_dDu\|X\|.
\]
This follows by expanding
$\widehat A X\widehat A^\dagger-A X A^\dagger$ into its
two linear error terms and its quadratic error term, using
$\|A\|\le1$, and summing at most $d$ terms.

Compute each transfer with fixed-point arithmetic of mesh $u$.
Compute one triangle of its matrix, set the other triangle to its
complex conjugate, and retain only the real part of each diagonal
entry. Every computed matrix is consequently Hermitian.
When its input has norm at most two, the additional arithmetic
error has operator norm at most $c_dD^3u$.
For example, directly evaluate the double sum for each entry:
there are at most $dD^2$ summands, each with a bounded number of
rounded scalar operations. Each entry error is at most
$c_dD^2u$ and matrix operator norm is at most $D$ times the
maximum entry absolute value. All partial sums have magnitude
at most $c_d(D+1)^2$, so $O_d(\log(D+1))$ integer guard bits
suffice. Conjugate completion and diagonal real-part selection
do not change this bound.

Let $e_j=\|\widehat R_j-R_j\|$.
Increasing a constant $L=c_d(D+1)^3$ absorbs the tensor and
arithmetic errors, and gives
\[
 e_{j-1}\le(1+Lu)e_j+Lu,\qquad e_N=0,
\]
as long as $e_j\le1$. Choose
\[
 b=\left\lceil
      N\log_2d+\log_2(16NL)\right\rceil.
\]
Then $NLu\le d^{-N}/16\le1/16$. Iterating the last recurrence
gives, for all $j$,
\[
 e_j\le(1+Lu)^N-1\le2NLu\le d^{-N}/8.
\]
The bound also justifies the assumed condition $e_j\le1$ at
every preceding step by induction; in particular, the computed
matrices have norm at most two.
This estimate holds simultaneously for every input word, every
prefix, and every candidate next output letter.

The resulting real scalar $\widehat p(z\mid y)$ is above
$d^{-N}/2$ exactly when $z$ is the correct prefix.
The comparison itself is an exact integer operation:
if $\widehat p=mu$ with signed integer $m$, compare
$2d^Nm$ with $2^b$.
Multiplication may use a constant multiple of the input bit
length for its intermediate integers; this does not change the
quadratic Boolean cost below. All integers used here have
$B=O_d(N+\log(N+1)+\log(D+1)+1)$ bits.

\subsection*{An exact reversible permutation}

Starting with the empty prefix, test the $d$ possible next
letters at each of the $N$ positions. Exactly one candidate
passes the threshold. This determines $\pi(y)$ exactly.
Recomputing every candidate probability from the full MPS uses
at most $dN^2$ transfer evaluations. The direct double-sum
method uses $O_d((D+1)^4)$ scalar operations per transfer.
Binary multiplication, rounding, addition, and the final
integer comparison use $O(B^2)$ Boolean operations per scalar
operation. Selection from the $d^2$ fixed tensor matrices per
site is a bounded-choice lookup and has polynomial cost within
the same estimate. Thus a Boolean circuit of size at most
$S=c_dN^2(D+1)^4B^2$ computes $\pi(y)$.
The rounded tensor table contains
$O_d(ND^2B)$ bits, also within this bound.

Retain the intermediate Boolean results, add the computed output
componentwise modulo $d$ into a second $N$-qudit register, and
reverse the calculation. This gives the exact unitary action
\[
 |y,z\rangle|0\rangle_{\rm work}
       \longmapsto
 |y,z+\pi(y)\rangle|0\rangle_{\rm work}.
\]
Each bounded-arity Boolean gate has an exact reversible
implementation with a constant number of gates on at most
three qudits. Boolean work values occupy the levels $0,1$;
the gates have a full unitary extension on the other levels.
For fixed $d$, an arbitrary three-qudit unitary has a
constant-size decomposition into arbitrary one- and two-qudit
gates. The exact synthesis includes its scalar phase.
Thus retaining and erasing the computation costs $O_d(S)$
gates and $O_d(S)$ initialized work qudits. The calculation is
correct for every input basis word, so its linear extension
applies to arbitrary coherent inputs.

The inverse permutation has the same operator cut ranks.
It also has bounded right-canonical tensors explicitly:
$B_j(a,b)=\overline{A_j(b,a)}$ represents the normalized
coefficients of $P_\pi^\dagger$ and satisfies the same canonical
identity. Apply the preceding construction to $\pi^{-1}$,
using subtraction in its output register. Starting with
$|y,0\rangle$, the two calculations give
\[
 |y,0\rangle\ \longmapsto\
 |y,\pi(y)\rangle\ \longmapsto\
 |y-\pi^{-1}(\pi(y)),\pi(y)\rangle
 =|0,\pi(y)\rangle.
\]
Swap the two $N$-qudit registers. The result is
$|\pi(y),0\rangle$, with all arithmetic work also at zero.
Every displayed operation has the specified coefficient $1$;
no input-dependent or scalar phase has been discarded.
This proves $CJ=JP_\pi$.

Finally, arrange the $O_d(S+N)=O_d(S)$ physical and work qudits
on a line. Route each two-qudit gate by swaps, apply it, and
undo the swaps. This multiplies the gate count by at most
$O_d(S)$ and leaves all initialized-output identities unchanged.
\end{proof}

\begin{corollary}\label{cor:joint-polynomial-qubit-permutation}
For $d=2$, the permutation circuit in
Theorem~\ref{thm:joint-polynomial-permutation} can be chosen
from the Clifford+$T$ gate alphabet, with the same polynomial
gate and clean auxiliary bounds, including on a line.
\end{corollary}

\begin{proof}
The arithmetic and comparisons above are exact Boolean
calculations. Their reversible implementations use only
NOT, CNOT, and Toffoli gates, including the bitwise additions
and subtractions that erase the input. All three gates have
exact Clifford+$T$ implementations with their specified scalar
phase; the phase-exact Toffoli construction is recorded in
\nolinkurl{docs/audits/2026-10-03_mpu_finite_gate_phase_lift.tex}.
Register swaps use three CNOT gates. Substituting these
constant-size implementations gives the claimed circuit.
\end{proof}

\begin{corollary}\label{cor:joint-polynomial-monomial}
Let $U$ be a monomial unitary on $N$ sites of fixed local
dimension $d$, with consecutive operator Schmidt ranks at
most $D$. For $0<\varepsilon<1$, it has a circuit with gate
count and auxiliary space polynomial jointly in
$N,D,\log(1/\varepsilon)$ whose exact initialized action is
$CJ=J\widetilde U$ for a unitary $\widetilde U$ satisfying
$\|\widetilde U-U\|_{\rm op}\le\varepsilon$.
All auxiliaries are returned exactly to zero.
\end{corollary}

\begin{proof}
Write $U|y\rangle=f(y)|\pi(y)\rangle$, with $|f(y)|=1$.
The permutation coefficients satisfy
$(P_\pi)_{a,y}=|U_{a,y}|^2$.
At any cut, the entrywise product of an operator Schmidt
expansion of $U$ with its complex conjugate gives a coefficient
factorization of $P_\pi$ with at most $D^2$ terms.
Thus its operator cut ranks are at most $D^2$, and the theorem
implements $P_\pi$ exactly with joint polynomial resources.

The operator $P_\pi^\dagger U$ is diagonal. Multiplying
operator Schmidt expansions shows that its cut ranks are at
most $D^3$. The diagonal approximation argument of
\nolinkurl{docs/audits/2026-10-02_mpu_diagonal_circuits.tex}
therefore gives a diagonal unitary $\widetilde D$ with
$\|\widetilde D-P_\pi^\dagger U\|_{\rm op}\le\varepsilon$
using joint polynomial resources and exact auxiliary erasure.
For local dimension $d$ the same proof normalizes its phase
vector by $d^{-N/2}$ and uses
$O_d(N+\log(ND+1)+\log(1/\varepsilon))$ precision bits;
the bounded-contraction and reversible phase-evaluation
steps are unchanged.
Compose this circuit with the exact permutation circuit.
Then $\widetilde U=P_\pi\widetilde D$ and left multiplication
by a unitary preserves operator norm, proving the conclusion.
\end{proof}

The workspace bound in these results is polynomial in $N,D$;
it is not the sharper $O_d(N\log(D+1))$ bound of the general
fixed-bond theorem. The permutation theorem is exact, whereas
the monomial conclusion approximates its arbitrary diagonal
phase factor. Neither result proves joint polynomial synthesis
for arbitrary MPUs that mix basis states coherently.
The inverse-erasure step uses the fact that the output is a
single basis word that determines the input basis word.
Preparing a general column $U|y\rangle$ while retaining $y$
does not provide this step: the map
$|y\rangle\mapsto|y\rangle U|y\rangle$ loses the inter-column
coherences after tracing out the retained label. Its physical
channel therefore differs from $\rho\mapsto U\rho U^\dagger$.
An auxiliary-only unitary cannot repair that difference.

\begin{thebibliography}{2}
\bibitem{exact-mpu-circuits}
TNLean contributors,
\emph{Exact polynomial circuits for matrix product unitaries},
2 October 2026, Section~5.
Local source:
\nolinkurl{docs/audits/2026-10-02_mpu_rank_two_circuits.tex}.
\bibitem{diagonal-mpu-circuits}
TNLean contributors,
\emph{Exact rank-two diagonal MPU circuits and approximation at arbitrary rank},
2 October 2026, Section~3.
Local source:
\nolinkurl{docs/audits/2026-10-02_mpu_diagonal_circuits.tex}.
\end{thebibliography}
\end{document}
